\section{Weekly Summary}

\begin{tcolorbox}[colback=yellow!10,colframe=black!50,title=AI-Generated Content]
\noindent The summary below was written by Claude (Anthropic) from that week's regression outputs and series descriptions. It has not been reviewed or fact-checked by a human, and should be read as an automated, informal recap -- not analysis.
\end{tcolorbox}

Another thrilling week of throwing unrelated economic data into a blender and pressing the button labeled "regress." As always, none of this means anything, and I would ask you not to phone the New Zealand government about any of it.

We kicked things off Monday with the World Pandemic Uncertainty Index for New Zealand versus, of all things, private construction spending on communication infrastructure in the United States. And folks, I regret to inform you that this pairing produced the statistical equivalent of a shrug. R-squared of less than one percent, a p-value hanging out around 0.32 like it forgot why it walked into the room. Even the trend mirage couldn't be bothered to show up. When we stripped the trends out things somehow got worse, which is genuinely impressive. Two series, zero relationship, both raw and detrended in complete agreement that they have never met and would like to keep it that way. A refreshingly honest null.

Tuesday is where the drama lives. We pitted a Northeast core CPI series against a nationwide transportation-services CPI series, and the raw regression came back with an R-squared of 0.97 and a p-value with more zeros in front of it than I care to count. "Astonishing," I said, "a genuine discovery," I did not say, because both of these are consumer price indexes that spend their entire lives marching upward together like two escalators bolted to the same wall. But here is the plot twist: even after we detrended and yanked that shared upward march out, the relationship held on with an R-squared of 0.84 and a p-value still buried in scientific notation. That is the rare, actually-kind-of-interesting case where two things move together even beyond their mutual habit of getting bigger over time. Two flavors of inflation turning out to be related is not exactly a Nobel-worthy shock, but by the standards of this project it's practically a soulmate match.

Then Wednesday arrived to restore cosmic balance. CPI for "Other Goods and Services" against a Medical Services Expenditures price index gave us a gorgeous raw R-squared of 0.98 — swoon — right up until we detrended it and the whole thing evaporated. R-squared collapsed to 0.00005 and the p-value strolled all the way up to 0.97, which in this business is the sound of a relationship being revealed as two things that merely aged in the same direction. Classic trend mirage, exactly the sort of thing this project exists to point at and laugh. So that's the week: one real-ish result, two whole lot of nothing, and the usual reminder that these are randomly paired series with no causal story intended. Spurious is the house style. See you next week for more.
